\documentclass[12pt]{article}
\usepackage[utf8]{inputenc}
\usepackage[T1]{fontenc}
\usepackage{hyperref}
\title{Gweno Language Resources for K-12 Curriculum Development}
\author{Compiled by Claude Sonnet 5.5 (claude-sonnet-5-5) for Andr\'as Kornai}
\date{}
\begin{document}
\maketitle

\section{Introduction}
Gweno (Kigweno, ISO 639-3 gwe) is a Bantu language of the North Pare Mountains, Kilimanjaro Region, Tanzania. Sources describe it as endangered and unwritten, spoken mostly by older adults, with younger people shifting to Asu and Swahili. Very little material exists.

\section{Primary Educational Resources}
None found. No primers or textbooks were identified in the searches.

\section{Linguistic and Reference Materials}
A Kigweno--Swahili--English lexicon (Sewangi 2008) is mentioned in the Wikipedia article; no direct link to it was seen in search results.
\begin{itemize}
\item Wikipedia, Gweno language: \url{https://en.wikipedia.org/wiki/Gweno_language}
\item Wikipedia, Ugweno (the region and its history): \url{https://en.wikipedia.org/wiki/Ugweno}
\end{itemize}

\section{Religious and Community Resources}
None found.

\section{Digital and Audio Resources}
The Endangered Languages Project lists a language sample page for Gweno: \url{https://www.endangeredlanguages.com/lang/4254/samples/9523?hl=de}

\section{Curriculum Development Context}
Gweno speakers are bilingual in Asu and mostly in Swahili, so any curriculum would need to be built from scratch with the community.
\end{document}